\section{Terminal states and sharp fixation time}\label{sec:temporal}
For a state $a$, write $m=\min_i a_i$ and $r=a-m\one$.
A positive run of $r$ is a maximal cyclic interval of positive coordinates.
Unless $r=0$, its zero coordinates separate these runs. The right endpoint
of a run is its last positive site in the chosen orientation.

\begin{lemma}\label{lem:minimum}
The minimum is invariant under $F$. A zero coordinate remains zero, and
\[
 F(a+c\one)=F(a)+c\one\qquad(c\in\Zge).
\]
\end{lemma}
\begin{proof}
Adding $c$ to every coordinate adds $c$ to every current, giving the
displayed identity. If $a_i=0$, then $q_{i-1}=q_i=0$, so $F(a)_i=0$.
The residual $r$ is nonnegative and has a zero, which persists. Therefore
$\min_i F(r)_i=0$, and translation gives $\min_i F(a)_i=m$.
\end{proof}

\begin{proposition}\label{prop:terminal}
Every orbit fixes. Its terminal residual has, at each original run's
original right endpoint, the total mass of that run, and is zero elsewhere.
The terminal state is obtained by adding $m$ back. In particular the
fixed states, and also all recurrent states, are precisely
\begin{equation}\label{eq:fixed}
 (a_i-m)(a_{i+1}-m)=0\qquad\text{for every }i.
\end{equation}
\end{proposition}
\begin{proof}
Work with $r$. Original zero boundaries persist and carry no current, so
each original run interval keeps its total mass $M$. If both $r_{i-1}$
and $r_i$ are positive, then
\[
 F(r)_i\ge\min(r_{i-1},r_i)>0.
\]
A positive site can therefore disappear only at the first site of its
current run. With its predecessor zero, it disappears exactly when its
mass is at most the next site's mass. Thus a run cannot split internally,
and can lose at most one positive site per update. It cannot acquire a
site across an existing zero.

Let $e$ be the original right endpoint, initially of mass $b>0$.
It never disappears: its next site remains zero, and hence
\[
 F(r)_e=r_e+\min(r_{e-1},r_e).
\]
While its run has at least two sites, this increases $r_e$ by at least
one. Since $r_e\le M$, there are at most $M-b$ nonfixed rounds in that
run. The run must reduce to the single mass $M$ at $e$, which is fixed.
Original runs are separated throughout, so the global fixation time is
the maximum of their fixation times.

If \eqref{eq:fixed} holds, all residual currents vanish. Conversely a
residual run of length at least two makes its endpoint increase, so the
state is not fixed. Uniform states correspond to $r=0$. Every orbit has
just been shown to fix; hence a recurrent state must itself be fixed.
\end{proof}

Define $\tau(a)$ to be the least $t\ge0$ such that $F^t(a)$ is fixed, and
put $H(n,N)=\max_{a\in X_{n,N}}\tau(a)$.

\begin{theorem}\label{thm:clock}
For the whole labelled carrier,
\begin{equation}\label{eq:height}
 H(n,N)=
 \begin{cases}
 0,&n\le2\text{ or }N=0,\\
 \lceil\log_2 N\rceil,&n=3,\ N\ge1,\\
 N-1,&n\ge4,\ N\ge1.
 \end{cases}
\end{equation}
\end{theorem}
\begin{proof}
For $n=1$ the only current cancels; for $n=2$ the two currents are equal.
Thus $F$ is the identity in both cases. The case $N=0$ is immediate.
For $n\ge3$, Proposition~\ref{prop:terminal} bounds every nontrivial run's
clock by $M-b\le N-1$.

For $n=3$, a nontrivial residual run has length two. Write its masses
as $(A,B)$, followed by zero, with $A,B>0$ and $M=A+B$. Its update is
\[
 (A,B)\longmapsto\bigl((A-B)_+,\,B+\min(A,B)\bigr).
\]
Induction gives, including all times after saturation,
\begin{equation}\label{eq:two-site}
 A_t=(M-2^tB)_+,\qquad B_t=\min(2^tB,M).
\end{equation}
Indeed $B_{t+1}=\min(2B_t,M)$ and $A_{t+1}=M-B_{t+1}$.
The exact clock is the least $t$ with $2^tB\ge M$, at most
$\lceil\log_2N\rceil$. For $N\ge2$, the state $(N-1,1,0)$ attains this
bound. For $N=1$, every state is fixed and the stated value is zero.

For $n\ge4$ and $N\ge3$, use $(N-2,1,1,0,\ldots,0)$. At times
$0\le t\le N-2$, its first three coordinates are
\[
 (N-2-t,\,1,\,1+t).
\]
While the first coordinate is positive, the middle coordinate receives
and sends one, and the last receives one. At time $N-2$ the state has
the pair $(1,N-1)$ at its last two occupied positions; one further update
gives the single mass $N$. Thus the clock is $N-1$.
For $N=2$, use $(1,1,0,\ldots,0)$, whose clock is one.
For $N=1$, all states are fixed. These witnesses attain every remaining
branch of \eqref{eq:height}.
\end{proof}

